\begin{afterword}
\summaryhead

The introduction to the second book, \textsc{How to Actually Change Your Mind}, opens with a 1951 football game between Dartmouth and Princeton. Asked who started the rough play, and later shown a film of the game, students from each school saw their own team as less to blame. When something we value is threatened, the introduction says, our perceptions and reasoning rally to defend it, and the stories we tell about our own reasons are often confabulated.

As a firmer standard it offers mathematics. Probability theory fixes the right degree of belief given priors and evidence. In the example, one of six classmates is a secret admirer, so the odds that it is Bob are 1:5; if people with a crush wink ten times as often and Bob winks, the odds become 2:1. Given modest constraints, the introduction concludes, every question of belief has an objectively right answer, and Robin Hanson is quoted: you are never entitled to your opinion. We will never be perfect Bayesians, but the ideal shows where we went wrong.

The second half describes the book's sequences, and ends by granting that biases are the substance of our reasoning. Debiasing is still possible, as arithmetic can be trained.

\responsehead

The introduction reports its sources with care, and its worked example is correct: 1:5 times 10:1 is 2:1, a probability of two thirds. The quotations from Hanson, Scott Alexander and Luke Muehlhauser match their posts. Its problems lie in the claims it builds on the example.

The central claim is that ``the question `What should I believe?' has an objectively right answer'' and that ``There is always a correct amount of confidence.'' The Bob example has a right answer because it fixes its own inputs: six equally likely candidates and a known ratio of ten. In a real case those inputs are themselves estimates. Consistency does not fix a prior, and whether anything else does divides Bayesians; subjective Bayesians hold that any coherent prior is permitted. The introduction states one side of this dispute as settled, and two paragraphs later drops even the condition ``Given our assumptions,'' calling the answer about Bob ``just as logically constrained'' as one in algebra. Our notes on ``How Much Evidence Does It Take?'' make the same point about starting probabilities.

The example is also loosely described at one step. Winking ``ten times as often'' becomes ``the 10:1 odds in favor of `a random person who winks at me has a crush on me'\,''. That is a likelihood ratio, which the introduction correctly names two paragraphs later. The odds that a random winker has a crush also depend on how many people have one. In a passage that teaches the method, the step reads as the base rate neglect that the first book's introduction describes.

The opening study is reported with a sharper contrast than its table shows. The 36\% of Dartmouth students who ``blamed Dartmouth'' are those who blamed Dartmouth alone. Another 53\% said both teams started it, so most Dartmouth students blamed their own team at least in part.

Finally, the introduction meets the doubt it raises. If biases are the ``very substance'' of our reasoning, can reading a book reduce them? Its answer is an analogy with arithmetic, which can be trained. As with ``Biases: An Introduction,'' no evidence is offered that the book's kind of teaching reduces bias.

In short: a careful worked example, from which the introduction draws a disputed conclusion, that every question of belief has an objectively right answer, and states it as settled.
\end{afterword}
